\chapter{2013年浙江卷（文）}

\section{单选题}

\begin{problem}
设集合 \( S = \{x \mid x > -2\} \)，\( T = \{x \mid -4 \leqslant x \leqslant 1\} \)，则 \( S \cap T = \)（\quad）

\choices
    {\([-4, +\infty)\)}
    {\((-2, +\infty)\)}
    {\([-4, 1]\)}
    {\((-2, 1]\)}
\end{problem}

\begin{answer}
D
\end{answer}

\begin{solution}
由交集的定义，既要满足 \(x>-2\)，又要满足 \(-4\leqslant x\leqslant1\)，所以
\(S\cap T=(-2,1]\)，故选 D.
\end{solution}

\begin{problem}
已知 \(\i\) 是虚数单位，则 \( (2 + \i)(3 + \i) = \)（\quad）

\choices
    {\(5 - 5\i\)}
    {\(7 - 5\i\)}
    {\(5 + 5\i\)}
    {\(7 + 5\i\)}

\end{problem}

\begin{answer}
C
\end{answer}

\begin{solution}
利用 \(\i^2=-1\)，展开得
\((2+\i)(3+\i)=6+2\i+3\i+\i^2=6+5\i-1=5+5\i\)，故选 C.
\end{solution}

\begin{problem}
若 \(\alpha \in \R\)，则“\(\alpha = 0\)”是“\(\sin \alpha < \cos \alpha\)”的（\quad）

\choices
    {充分不必要条件}
    {必要不充分条件}
    {充分必要条件}
    {既不充分也不必要条件}

\end{problem}

\begin{answer}
A
\end{answer}

\begin{solution}
当 \(\alpha=0\) 时，\(\sin\alpha=0<1=\cos\alpha\)，所以“\(\alpha=0\)”是“\(\sin\alpha<\cos\alpha\)”的充分条件. 又例如 \(\alpha=-\frac{\pi}{4}\) 时，\(\sin\alpha<\cos\alpha\)，但 \(\alpha\ne0\)，所以不是必要条件，故选 A.
\end{solution}

\begin{problem}
设 \(m\)，\(n\) 是两条不同的直线，\(\alpha\)，\(\beta\) 是两个不同的平面（\quad）

\choices
    {若 \(m\parallel\alpha\)，\(n\parallel\alpha\)，则 \(m\parallel n\)}
    {若 \(m\parallel\alpha\)，\(m\parallel\beta\)，则 \(\alpha\parallel\beta\)}
    {若 \(m\parallel n\)，\(m \perp \alpha\)，则 \(n \perp \alpha\)}
    {若 \(m\parallel\alpha\)，\(\alpha \perp \beta\)，则 \(m \perp \beta\)}

\end{problem}

\begin{answer}
C
\end{answer}

\begin{solution}
选项 C 中，\(m\parallel n\)，且 \(m\perp\alpha\). 直线 \(m\) 垂直于平面 \(\alpha\) 时，所有与 \(m\) 平行的直线都垂直于 \(\alpha\)，所以 \(n\perp\alpha\)，C 正确. 其余选项均不成立：对 A，取平面 \(\alpha\) 为 \(z=0\)，取同在平面 \(z=1\) 内且相交的两条直线 \(m\)、\(n\)，则它们都平行于 \(\alpha\)，但不平行；对 B，取 \(\alpha\colon z=0\)、\(\beta\colon y=0\)，取平行于 \(x\) 轴且位于 \(y=z=1\) 的直线 \(m\)，则 \(m\parallel\alpha\)、\(m\parallel\beta\)，而 \(\alpha\) 与 \(\beta\) 相交；对 D，取 \(\alpha\colon z=0\)、\(\beta\colon x=0\)，取平行于 \(y\) 轴且位于 \(x=z=1\) 的直线 \(m\)，则 \(\alpha\perp\beta\)、\(m\parallel\alpha\)，但 \(m\) 不垂直于 \(\beta\)，故选 C.
\end{solution}

\begin{problem}
已知某几何体的三视图（单位：\(\mathrm{cm}\)）如图所示，则该几何体的体积是（\quad）

\bitmapfigure[max width=0.62\linewidth]{img/2013/zhejiang_liberal/q5_fig1.png}

\choices
    {\(108\mathrm{~cm}^3\)}
    {\(100\mathrm{~cm}^3\)}
    {\(92\mathrm{~cm}^3\)}
    {\(84\mathrm{~cm}^3\)}

\end{problem}

\begin{answer}
B
\end{answer}

\begin{solution}
由三视图可还原为一个长、宽、高分别为 \(6\mathrm{~cm}\)、\(6\mathrm{~cm}\)、\(3\mathrm{~cm}\) 的长方体，并去掉一个以三条互相垂直的棱为棱、棱长分别为 \(4\mathrm{~cm}\)、\(4\mathrm{~cm}\)、\(3\mathrm{~cm}\) 的三棱锥. 因此
\[
V=6\times6\times3-\frac{1}{3}\times\frac{1}{2}\times4\times4\times3=108-8=100\mathrm{~cm}^3
\]
故选 B.
\end{solution}

\begin{problem}
函数 \( f(x) = \sin x\cos x + \frac{\sqrt{3}}{2}\cos 2x \) 的最小正周期和振幅分别是（\quad）

\choices
    {\(\pi\)，\(1\)}
    {\(\pi\)，\(2\)}
    {\(2\pi\)，\(1\)}
    {\(2\pi\)，\(2\)}

\end{problem}

\begin{answer}
A
\end{answer}

\begin{solution}
由二倍角公式，
\[
f(x)=\frac{1}{2}\sin2x+\frac{\sqrt3}{2}\cos2x=\sin\left(2x+\frac{\pi}{3}\right)
\]
因此角变量的系数为 \(2\)，最小正周期为 \(T=\frac{2\pi}{2}=\pi\)，且正弦函数的振幅为 \(1\)，故选 A.
\end{solution}

\begin{problem}
已知 \(a\)，\(b\)，\(c \in \R\)，函数 \(f(x) = ax^2 + bx + c\)，若 \(f(0) = f(4) > f(1)\)，则（\quad）

\choices
    {\(a > 0\)，\(4a + b = 0\)}
    {\(a < 0\)，\(4a + b = 0\)}
    {\(a > 0\)，\(2a + b = 0\)}
    {\(a < 0\)，\(2a + b = 0\)}

\end{problem}

\begin{answer}
A
\end{answer}

\begin{solution}
由 \(f(0)=f(4)\) 得
\[
c=16a+4b+c
\]
所以 \(4a+b=0\). 又由 \(f(0)>f(1)\) 得
\[
c>a+b+c
\]
即 \(a+b<0\). 代入 \(b=-4a\)，得 \(-3a<0\)，所以 \(a>0\)，故选 A.
\end{solution}

\begin{problem}
已知函数 \( y = f(x) \) 的图象是下列四个图象之一，且其导函数 \( y = f'(x) \) 的图象如图所示，则该函数的图象是（\quad）

\FigureLayoutDeclare{img/2013/zhejiang_liberal/q8_fig1.png}{5.6cm}{11bp 14bp 6bp 0bp}
\bitmapfigure[max width=6.0cm]{img/2013/zhejiang_liberal/q8_fig1.png}

\choices
    {\choicebitmap[width=0.15\paperwidth]{img/2013/zhejiang_liberal/q8_choice_a.png}}
    {\choicebitmap[width=0.15\paperwidth]{img/2013/zhejiang_liberal/q8_choice_b.png}}
    {\choicebitmap[width=0.15\paperwidth]{img/2013/zhejiang_liberal/q8_choice_c.png}}
    {\choicebitmap[width=0.15\paperwidth]{img/2013/zhejiang_liberal/q8_choice_d.png}}

\end{problem}

\begin{answer}
B
\end{answer}

\begin{solution}
由图可知，在 \((-1,1)\) 内 \(f'(x)>0\)，所以 \(f(x)\) 在该区间内递增. 又因为 \(f'(x)\) 在 \((-1,0)\) 上递增、在 \((0,1)\) 上递减，所以 \(f(x)\) 在 \((-1,0)\) 上为凸函数，在 \((0,1)\) 上为凹函数，且在 \(x=0\) 处切线斜率达到最大值. 四个图象中只有 B 同时满足这些单调性、凹凸性和斜率变化，故选 B.
\end{solution}

\begin{problem}
如图，\(F_{1}\)，\(F_{2}\) 是椭圆 \(C_{1}: \frac{x^{2}}{4} + y^{2} = 1\) 与双曲线 \(C_{2}\) 的公共焦点，\(A\)，\(B\) 分别是 \(C_{1}\)，\(C_{2}\) 在第二、四象限的公共点. 若四边形 \(AF_{1}BF_{2}\) 为矩形，则 \(C_{2}\) 的离心率是（\quad）

\FigureLayoutDeclare{img/2013/zhejiang_liberal/q9_fig1.png}{7.2cm}{199bp 0bp 91bp 0bp}
\bitmapfigure[max width=0.58\linewidth]{img/2013/zhejiang_liberal/q9_fig1.png}

\choices
    {\(\sqrt{2}\)}
    {\(\sqrt{3}\)}
    {\( \frac{3}{2} \)}
    {\(\frac{\sqrt{6}}{2}\)}

\end{problem}

\begin{answer}
D
\end{answer}

\begin{solution}
椭圆 \(C_1\) 的长半轴为 \(2\)，半焦距为
\(c=\sqrt{2^2-1^2}=\sqrt3\)，故 \(|F_1F_2|=2c=2\sqrt3\). 令
\(x=|AF_1|\)，\(y=|AF_2|\). 由椭圆定义，
\[
x+y=4
\]
又因四边形 \(AF_1BF_2\) 为矩形，\(\angle F_1AF_2=90^\circ\)，在直角三角形 \(AF_1F_2\) 中
\[
x^2+y^2=|F_1F_2|^2=12
\]
因为点 \(A\) 在第二象限，设其横坐标为 \(x_A<0\)，则
\[
|AF_2|^2-|AF_1|^2=(x_A-\sqrt3)^2-(x_A+\sqrt3)^2=-4\sqrt3 x_A>0
\]
所以 \(AF_1<AF_2\)，从而 \(x=2-\sqrt2\)、\(y=2+\sqrt2\). 双曲线 \(C_2\) 的焦点为 \(F_1\)，\(F_2\)，且点 \(A\) 在左支上，所以
\(2a=y-x=2\sqrt2\)，\(2c=2\sqrt3\).
因此其离心率为
\[
e=\frac{c}{a}=\frac{\sqrt3}{\sqrt2}=\frac{\sqrt6}{2}
\]
故选 D.
\end{solution}

\begin{problem}
设 \(a\)，\(b \in \R\)，定义运算“\(\wedge\)”和“\(\vee\)”如下：
\(a \wedge b = \begin{cases}
a, & a \leqslant b,\\
b, & a > b,
\end{cases}\)
\(a \vee b = \begin{cases}
b, & a \leqslant b,\\
a, & a > b.
\end{cases}\) 若正数 \(a\)，\(b\)，\(c\)，\(d\) 满足 \(ab \geqslant 4\)，\(c + d \leqslant 4\)，则（\quad）

\choices
    {\(a \wedge b \geqslant 2\)，\(c \wedge d \leqslant 2\)}
    {\(a \wedge b \geqslant 2\)，\(c \vee d \geqslant 2\)}
    {\(a \vee b \geqslant 2\)，\(c \wedge d \leqslant 2\)}
    {\(a \vee b \geqslant 2\)，\(c \vee d \geqslant 2\)}

\end{problem}

\begin{answer}
C
\end{answer}

\begin{solution}
由定义，\(a\wedge b=\min\{a,b\}\)，\(a\vee b=\max\{a,b\}\). 若 \(a\vee b<2\)，则 \(a<2\)、\(b<2\)，从而 \(ab<4\)，与 \(ab\geqslant4\) 矛盾，所以 \(a\vee b\geqslant2\). 若 \(c\wedge d>2\)，则 \(c>2\)、\(d>2\)，从而 \(c+d>4\)，与 \(c+d\leqslant4\) 矛盾，所以 \(c\wedge d\leqslant2\). 因而选项 C 的两个不等式恒成立. 选项 A、B 均可由 \(a=1\)、\(b=4\)、\(c=d=1\) 否定，因为此时 \(ab=4\)、\(c+d=2\)，但 \(a\wedge b=1<2\)；选项 D 可由 \(a=b=2\)、\(c=d=1\) 否定，因为此时 \(c\vee d=1<2\). 故选 C.
\end{solution}

\section{填空题}

\begin{problem}
已知函数 \( f(x) = \sqrt{x - 1} \). 若 \( f(a) = 3 \)，则实数 \( a = \)\fillinblank{}.
\end{problem}

\begin{answer}
10
\end{answer}

\begin{solution}
由 \(f(a)=3\) 得 \(\sqrt{a-1}=3\). 函数定义域要求 \(a\geqslant1\)，两边平方得 \(a-1=9\)，所以 \(a=10\).
\end{solution}

\begin{problem}
从 \(3\) 男 \(3\) 女共 \(6\) 名同学中任选 \(2\) 名（每名同学被选中的机会均等）. 这 \(2\) 名都是女同学的概率等于\fillinblank{}.
\end{problem}

\begin{answer}
\(\frac{1}{5}\)
\end{answer}

\begin{solution}
从 \(6\) 名同学中任选 \(2\) 名共有 \(\mathrm C_6^2=15\) 种等可能结果，其中 \(2\) 名都是女生的结果有 \(\mathrm C_3^2=3\) 种. 因此所求概率为
\[
P=\frac{\mathrm C_3^2}{\mathrm C_6^2}=\frac{3}{15}=\frac{1}{5}
\]
\end{solution}

\begin{problem}
直线 \( y = 2x + 3 \) 被圆 \( x^{2} + y^{2} - 6x - 8y = 0 \) 所截得的弦长等于\fillinblank{}.
\end{problem}

\begin{answer}
\(4\sqrt{5}\)
\end{answer}

\begin{solution}
将圆方程配方得
\[
(x-3)^2+(y-4)^2=25
\]
所以圆心为 \(O(3,4)\)，半径为 \(5\). 直线方程化为 \(2x-y+3=0\)，圆心到直线的距离为
\[
d=\frac{|2\times3-4+3|}{\sqrt{2^2+(-1)^2}}=\sqrt5
\]
弦长的一半为 \(\sqrt{5^2-d^2}=\sqrt{25-5}=2\sqrt5\)，故弦长为 \(4\sqrt5\).
\end{solution}

\begin{problem}
若某程序框图如图所示，则该程序运行后输出的值等于\fillinblank{}.

\texfigure[max width=0.46\linewidth]{img_repaint/2013/zhejiang_liberal/q14_fig1.tex}
\end{problem}

\begin{answer}
\(\frac{9}{5}\)
\end{answer}

\begin{solution}
程序从 \(S=1\)、\(k=1\) 开始，在 \(k\leqslant4\) 时依次加上 \(\frac{1}{k(k+1)}\)，然后令 \(k\) 增加 1. 因此输出值为
\[
S=1+\sum_{k=1}^{4}\frac{1}{k(k+1)}
=1+\sum_{k=1}^{4}\left(\frac{1}{k}-\frac{1}{k+1}\right)
=1+1-\frac{1}{5}=\frac{9}{5}
\]
\end{solution}

\begin{problem}
设 \(z = kx + y\)，其中实数 \(x\)，\(y\) 满足 \(\begin{cases}
x \geqslant 2,\\
x - 2y + 4 \geqslant 0,\\
2x - y - 4 \leqslant 0.
\end{cases}\) 若 \(z\) 的最大值为 12，则实数 \(k = \)\fillinblank{}.
\end{problem}

\begin{answer}
2
\end{answer}

\begin{solution}
约束条件化为
\(y\leqslant\frac{x}{2}+2\)，\(y\geqslant2x-4\)，\(x\geqslant2\).
由两条斜边界相交得 \(x=4\)，\(y=4\)，而在直线 \(x=2\) 上两条边界分别给出 \(y=0\) 和 \(y=3\). 又因可行域要求 \(2x-4\leqslant\frac{x}{2}+2\)，所以 \(2\leqslant x\leqslant4\)，可行域的顶点为 \((2,0)\)、\((2,3)\)、\((4,4)\). 线性函数在三角形可行域上的最大值取在顶点处，在三个顶点处分别有
\(z=2k\)，\(z=2k+3\)，\(z=4k+4\).
点 \((2,0)\) 的值小于点 \((2,3)\) 的值. 若最大值在 \((2,0)\) 处达到 12，则 \(k=6\)，但此时 \((2,3)\) 处的值为 15；若最大值在 \((2,3)\) 处达到 12，则 \(k=\frac{9}{2}\)，但此时 \((4,4)\) 处的值为 22. 因此最大值只能在 \((4,4)\) 处达到，满足 \(4k+4=12\)，从而 \(k=2\). 此时三个顶点处的函数值为 \(4\)、\(7\)、\(12\)，确实最大值为 12.
\end{solution}

\begin{problem}
设 \(a\)，\(b \in \R\)，若 \(x \geqslant 0\) 时恒有 \(0 \leqslant x^4 - x^3 + ax + b \leqslant (x^2 - 1)^2\)，则 \(ab = \)\fillinblank{}.
\end{problem}

\begin{answer}
\(-1\)
\end{answer}

\begin{solution}
令 \(P(x)=x^4-x^3+ax+b\). 取 \(x=1\)，由
\(0\leqslant P(1)\leqslant(1-1)^2=0\)，得
\[
a+b=0
\]
于是 \(b=-a\)，并且
\[
P(x)=x^4-x^3+a(x-1)=(x-1)(x^3+a)
\]
当 \(0\leqslant x<1\) 时，\(x-1<0\)，为使 \(P(x)\geqslant0\)，必须有 \(x^3+a\leqslant0\)；令 \(x\to1^-\) 得 \(a+1\leqslant0\). 当 \(x>1\) 时，必须有 \(x^3+a\geqslant0\)；令 \(x\to1^+\) 得 \(a+1\geqslant0\). 所以 \(a=-1\)，\(b=1\). 此时
\(P(x)=(x-1)^2(x^2+x+1)\)，\((x^2-1)^2-P(x)=x(x-1)^2\)
对所有 \(x\geqslant0\)，有 \(x^2+x+1>0\)、\(x(x-1)^2\geqslant0\)，所以确实满足题设不等式，故 \(ab=-1\).
\end{solution}

\begin{problem}
设 \(\bs{e}_1\)，\(\bs{e}_2\) 为单位向量，非零向量 \(\bs{b} = x\bs{e}_1 + y\bs{e}_2\)，\(x\)，\(y \in \R\). 若 \(\bs{e}_1\)，\(\bs{e}_2\) 的夹角为 \(\frac{\pi}{6}\)，则 \(\frac{|x|}{|\bs{b}|}\) 的最大值等于\fillinblank{}.
\end{problem}

\begin{answer}
2
\end{answer}

\begin{solution}
若 \(x=0\)，则 \(\frac{|x|}{|\bs{b}|}=0\). 设 \(x\ne0\)，令 \(t=\frac{y}{x}\). 由两个单位向量夹角为 \(\frac{\pi}{6}\)，有
\[
|\bs{b}|^2=x^2+y^2+2xy\cos\frac{\pi}{6}=x^2+\sqrt3xy+y^2
\]
因此
\[
\frac{|x|}{|\bs{b}|}
=\frac{1}{\sqrt{1+\sqrt3t+t^2}}
=\frac{1}{\sqrt{\left(t+\frac{\sqrt3}{2}\right)^2+\frac14}}
\leqslant2
\]
当 \(t=-\frac{\sqrt3}{2}\) 时等号成立，且 \(\bs{b}\ne\bs{0}\)，所以最大值为 \(2\).
\end{solution}

\section{解答题}

\begin{problem}
在锐角 \(\triangle ABC\) 中，内角 \(A\)，\(B\)，\(C\) 的对边分别为 \(a\)，\(b\)，\(c\)，且 \(2a \sin B = \sqrt{3} b\).
\begin{enumerate}
\item 求角 \(A\) 的大小；
\item 若 \(a = 6\)，\(b + c = 8\)，求 \(\triangle ABC\) 的面积.
\end{enumerate}
\end{problem}

\begin{solution}
\begin{enumerate}
\item 由正弦定理 \(\frac{a}{\sin A}=\frac{b}{\sin B}\)，原条件化为
\[
2\sin A\sin B=\sqrt3\sin B
\]
三角形为锐角三角形，故 \(\sin B\ne0\)，从而 \(\sin A=\frac{\sqrt3}{2}\). 又 \(A\) 为锐角，所以 \(A=\frac{\pi}{3}\).

\item 由余弦定理和 \(\cos A=\frac12\)，
\[
36=b^2+c^2-2bc\cos A=b^2+c^2-bc=(b+c)^2-3bc=64-3bc
\]
所以 \(bc=\frac{28}{3}\). 于是
\[
S_{\triangle ABC}=\frac12bc\sin A
=\frac12\times\frac{28}{3}\times\frac{\sqrt3}{2}
=\frac{7\sqrt3}{3}
\]
\end{enumerate}
\end{solution}

\begin{problem}
在公差为 \(d\) 的等差数列 \(\{a_{n}\}\) 中，已知 \(a_1 = 10\)，且 \(a_1\)，\(2a_2 + 2\)，\(5a_3\) 成等比数列.
\begin{enumerate}
\item 求 \(d\)，\(a_{n}\)；
\item 若 \(d < 0\)，求 \( \left|a_{1}\right| + \left|a_{2}\right| + \left|a_{3}\right| + \cdots + \left|a_{n}\right| \).
\end{enumerate}
\end{problem}

\begin{solution}
\begin{enumerate}
\item 等差数列的通项为 \(a_n=10+(n-1)d\)，故 \(a_2=10+d\)，\(a_3=10+2d\). 由三项成等比数列，
\[
(2a_2+2)^2=a_1\cdot5a_3
\]
即
\[
(22+2d)^2=10(50+10d)
\]
整理得 \(d^2-3d-4=0\)，所以 \(d=-1\) 或 \(d=4\). 相应地
\(a_n=10+(n-1)(-1)=11-n\)，\(a_n=10+(n-1)4=4n+6\).

\item 因为 \(d<0\)，只能取 \(d=-1\)，此时 \(a_n=11-n\). 设 \(n\in\Z_{>0}\). 当 \(n\leqslant11\) 时，前 \(n\) 项均非负，故
\[
\sum_{k=1}^{n}|a_k|=\sum_{k=1}^{n}(11-k)=\frac{21n-n^2}{2}
\]
当 \(n\geqslant12\) 时，前 11 项的绝对值和为 \(10+9+\cdots+0=55\)，且
\[
\sum_{k=12}^{n}|a_k|=\sum_{k=12}^{n}(k-11)=\frac{(n-11)(n-10)}{2}
\]
所以前 \(n\) 项的绝对值和等于 \(55+\frac{(n-11)(n-10)}{2}\).
所以
\[
\sum_{k=1}^{n}|a_k|=
\begin{cases}
-\frac{1}{2}n^2+\frac{21}{2}n, & n\leqslant11,\\
\frac{1}{2}n^2-\frac{21}{2}n+110, & n\geqslant12.
\end{cases}
\]
\end{enumerate}
\end{solution}

\begin{problem}
如图，在四棱锥 \(P - ABCD\) 中，\(PA \perp\) 平面 \(ABCD\)，\(AB = BC = 2\)，\(AD = CD = \sqrt{7}\)，\(PA = \sqrt{3}\)，\(\angle ABC = 120^\circ\)，\(G\) 为线段 \(PC\) 上的点.
\begin{enumerate}
\item 证明：\(BD \perp\) 平面 \(APC\)；
\item 若 \(G\) 为 \(PC\) 的中点，求 \(DG\) 与平面 \(APC\) 所成的角的正切值；
\item 若 \(G\) 满足 \(PC \perp \) 平面 \(BGD\)，求 \(\frac{PG}{GC}\) 的值.
\end{enumerate}

\bitmapfigure[max width=0.40\linewidth]{img/2013/zhejiang_liberal/q20_fig1.png}
\end{problem}

\begin{solution}
\begin{enumerate}
\item 因为 \(AB=BC\)，点 \(B\) 在线段 \(AC\) 的垂直平分线上；又因为 \(AD=CD\)，点 \(D\) 也在线段 \(AC\) 的垂直平分线上. 因此 \(BD\perp AC\). 又 \(PA\perp\) 平面 \(ABCD\)，所以 \(PA\perp BD\). 直线 \(AC\)、\(PA\) 是平面 \(APC\) 内相交的两条直线，故 \(BD\perp\) 平面 \(APC\).

\item 设 \(O=AC\cap BD\). 在 \(\triangle ABC\) 中，由余弦定理
\[
AC^2=AB^2+BC^2-2AB\cdot BC\cos120^\circ=4+4+4=12
\]
所以 \(AC=2\sqrt3\)，\(OC=\sqrt3\). 在直角三角形 \(COD\) 中，
\[
OD=\sqrt{CD^2-OC^2}=\sqrt{7-3}=2
\]
当 \(G\) 为 \(PC\) 的中点时，在 \(\triangle PAC\) 中，\(O\)、\(G\) 分别为 \(AC\)、\(PC\) 的中点，故
\(OG\parallel PA\)，\(OG=\frac12PA=\frac{\sqrt3}{2}\).
从而 \(OG\) 在平面 \(APC\) 内，且 \(DO\perp\) 平面 \(APC\)，所以 \(\angle DGO\) 是 \(DG\) 与平面 \(APC\) 所成的角. 在直角三角形 \(DOG\) 中，
\[
\tan\angle DGO=\frac{OD}{OG}=\frac{2}{\sqrt3/2}=\frac{4\sqrt3}{3}
\]

\item 由 \(PC\perp\) 平面 \(BGD\)，且 \(O\)、\(G\) 都在平面 \(BGD\) 内，所以 \(PC\perp OG\). 又
\[
PC=\sqrt{PA^2+AC^2}=\sqrt{3+12}=\sqrt{15}
\]
因为 \(PC\perp OG\)，所以 \(\angle OGC=90^\circ\)；又因 \(PA\perp AC\)，所以 \(\angle CAP=90^\circ\). 另外，\(OC\) 与 \(AC\) 共线、\(CG\) 与 \(CP\) 共线，故 \(\angle OCG=\angle PCA\). 因此
\[
\triangle COG\sim\triangle CPA
\]
对应边给出
\[
\frac{GC}{AC}=\frac{OC}{PC}
\]
因此
\[
GC=AC\cdot\frac{OC}{PC}=2\sqrt3\cdot\frac{\sqrt3}{\sqrt{15}}=\frac{2\sqrt{15}}{5}
\]
于是
\(PG=PC-GC=\sqrt{15}-\frac{2\sqrt{15}}{5}=\frac{3\sqrt{15}}{5}\)，\(\frac{PG}{GC}=\frac{3\sqrt{15}/5}{2\sqrt{15}/5}=\frac32\).
\end{enumerate}
\end{solution}

\begin{problem}
已知 \(a \in \R\)，函数 \(f(x) = 2x^3 - 3(a + 1)x^2 + 6ax\).
\begin{enumerate}
\item 若 \(a = 1\)，求曲线 \(y = f(x)\) 在点 \((2, f(2))\) 处的切线方程；
\item 若 \( |a| > 1 \)，求 \( f(x) \) 在闭区间 \([0, 2|a|]\) 上的最小值.
\end{enumerate}
\end{problem}

\begin{solution}
\begin{enumerate}
\item 当 \(a=1\) 时，\(f(x)=2x^3-6x^2+6x\)，故
\(f(2)=4\)，\(f'(x)=6x^2-12x+6\)，\(f'(2)=6\).
所以切线方程为
\[
y-4=6(x-2)
\]
即 \(y=6x-8\).

\item 对任意 \(a\)，
\[
f'(x)=6x^2-6(a+1)x+6a=6(x-1)(x-a)
\]
若 \(a<-1\)，区间为 \([0,-2a]\)，其中只有临界点 \(x=1\). 函数在 \([0,1]\) 上递减、在 \([1,-2a]\) 上递增，故最小值为
\[
f(1)=2-3(a+1)+6a=3a-1
\]

若 \(a>1\)，区间为 \([0,2a]\)，临界点为 \(x=1\)，\(a\). 函数在 \([0,1]\) 上递增、在 \([1,a]\) 上递减、在 \([a,2a]\) 上递增. 计算端点和极小点的函数值：
\(f(0)=0\)，\(f(a)=a^2(3-a)\)，\(f(2a)=4a^3>0\).
因此只需比较 \(0\) 与 \(a^2(3-a)\)：当 \(1<a\leqslant3\) 时最小值为 \(0\)；当 \(a>3\) 时最小值为 \(a^2(3-a)\). 综上，最小值为
\[
\begin{cases}
3a-1, & a<-1,\\
0, & 1<a\leqslant3,\\
a^2(3-a), & a>3.
\end{cases}
\]
\end{enumerate}
\end{solution}

\begin{problem}
已知抛物线 \(C\) 的顶点为 \(O(0,0)\). 焦点为 \(F(0,1)\).
\begin{enumerate}
\item 求抛物线 \(C\) 的方程；
\item 过点 \(F\) 作直线交抛物线 \(C\) 于 \(A\)，\(B\) 两点. 若直线 \(AO\)，\(BO\) 分别交直线 \(l: y = x - 2\) 于 \(M\)，\(N\) 两点，求 \(|MN|\) 的最小值.
\end{enumerate}
\end{problem}

\begin{solution}
\begin{enumerate}
\item 设抛物线方程为 \(x^2=2py\)，其中 \(p>0\). 其焦点为 \((0,\frac{p}{2})\)，由焦点 \(F(0,1)\) 得 \(p=2\)，所以抛物线方程为 \(x^2=4y\).

\item 过焦点的竖直直线为 \(x=0\)，只与抛物线在顶点 \(O\) 处相切，不能交于两个不同的点 \(A\)，\(B\). 因此过焦点且能交抛物线于两点的直线可设为
\[
y=kx+1
\]
设其与抛物线交于 \(A(x_1,y_1)\)、\(B(x_2,y_2)\)，联立得
\[
x^2-4kx-4=0
\]
所以
\(x_1+x_2=4k\)，\(x_1x_2=-4\)，\(|x_1-x_2|=4\sqrt{k^2+1}\).
由于 \(y_i=\frac{x_i^2}{4}\)，直线 \(AO\) 的方程为 \(y=\frac{x_1}{4}x\). 与 \(l:y=x-2\) 联立，得点 \(M\) 的横坐标
\[
x_M=\frac{8}{4-x_1}
\]
同理，
\[
x_N=\frac{8}{4-x_2}
\]
直线 \(l\) 的斜率为 1，因此
\[
|MN|=\sqrt2|x_M-x_N|=\sqrt2\left|\frac{8}{4-x_1}-\frac{8}{4-x_2}\right|=\frac{8\sqrt{2(k^2+1)}}{|4k-3|}
\]
这里 \(4k-3\ne0\)，否则相应的交点不为有限点. 由
\[
25(k^2+1)-(4k-3)^2=(3k+4)^2\geqslant0
\]
可得
\[
\frac{\sqrt{k^2+1}}{|4k-3|}\geqslant\frac15
\]
当 \(k=-\frac43\) 时等号成立，所以 \(|MN|\) 的最小值为
\[
\frac{8\sqrt2}{5}
\]
\end{enumerate}
\end{solution}
